\section{Журнал выполнения}
\begin{enumerate}
    \item Заголовков и ссылок нет в текстовых файлах ЛР1, они лежат в \texttt{meta.jsonl} (JSON). Чтобы не писать разбор JSON на C++, список документов с заголовками и URL готовит скрипт на Python (обвязка), а индексатор читает простой формат \enquote{путь, заголовок, URL через табуляцию}.
    \item По правилам курса структуры индекса нельзя брать из STL, поэтому динамический массив, хеш-таблица и сортировка написаны самостоятельно. Динамический массив копирует элементы через \texttt{memcpy}/\texttt{realloc}, поэтому записи словаря --- простые структуры, а список документов терма --- указатель и ёмкость внутри записи.
    \item Проверка полноты: число термов в индексе (3\,974\,421) совпало с числом различных терминов в частотном списке ЛР5, построенном тем же токенизатором независимым скриптом на Python.
    \item Длина строки в пуле записывается двумя байтами, поэтому токен длиннее 65\,535 байт не поместился бы. Такие токены индексатор пропускает и считает (\texttt{skipped\_tokens}); в корпусе их нет, самый длинный токен --- 754 символа (кусок ссылки в base64 из поста Reddit).
    \item Время индексации одного и того же корпуса менялось от запуска к запуску до 2 раз: процессор гибридный (см. ЛР3), на машине работают другие процессы, а на корпусе в 1 млн документов главное --- сколько файлов осталось в кэше ОС после предыдущего запуска (в ЛР11 один и тот же прогон занимал от 152 до 284~с). Каждый замер повторяется 3 раза, берётся минимальное время.
    \item Корпус заменён на 1\,000\,000 документов четырёх источников (ЛР1). Скрипт \texttt{doclist.py} перебирал только каталоги \texttt{wikivoyage} и \texttt{wikipedia}; теперь он сам находит источники в каталоге корпуса (сначала Wikivoyage и Wikipedia, затем остальные по алфавиту). Индексатор не менялся: индекс миллиона строится в памяти за 2,5~мин и занимает 2,1~ГиБ памяти.
\end{enumerate}

\section{План тестирования}
Тесты --- юнит-тесты на C++ (\texttt{tests/test\_index.cpp}) на собственном минимальном каркасе \texttt{lab3/tests/unittest.h} (макросы \texttt{TEST} и \texttt{CHECK}, без сторонних библиотек). Они вызывают функции программы напрямую, а последний тест запускает утилиты командной строки как отдельные процессы. Сборка и запуск --- \texttt{make test}. Эталонные значения в тестах считаются независимо от проверяемого кода, для этого используются контейнеры STL (\texttt{std::map}, \texttt{std::set}) --- это код тестов, а не поисковой системы. Файл индекса тест разбирает собственным читателем, написанным строго по описанию формата (раздел \enquote{Формат индекса}), --- так проверяется и индексатор, и само описание.

\begin{center}
\small
\begin{tabular}{|c|p{0.3\textwidth}|p{0.58\textwidth}|}
\hline
\textnumero & Тест & Что проверяется \\
\hline
1 & Маленький корпус & Словарь и списки совпадают с эталоном, термы упорядочены; поля заголовка (сигнатура, версия, размеры, $N$, $T$, $P$); разделы идут подряд; статистика сборщика совпадает с файлом. \\
\hline
2 & Регистр, Юникод, повторы & \texttt{Paris PARIS paris}, \texttt{Москва МОСКВА} --- по одному терму; документ в списке один раз; порядок термов по байтам UTF-8. \\
\hline
3 & Прямой индекс & Заголовки и URL в UTF-8, число токенов (в т.\,ч. пустой документ); читатель индекса \texttt{index\_reader.h} выдаёт те же данные. \\
\hline
4 & Двоичный поиск терма & \texttt{Index::find}: существующие термы с правильным df и списком, первый терм таблицы; отсутствующий терм, его префикс и продолжение --- $-1$. \\
\hline
5 & Случайный корпус & 400 документов, около 13 тыс. термов: хеш-таблица трижды перестраивается, длинные списки растут; весь индекс совпадает с эталоном, \texttt{find} находит каждый терм. \\
\hline
6 & Строка списка документов & Разбор \enquote{путь, заголовок, URL}; строка без заголовка или URL --- ошибка. \\
\hline
7 & Пустой корпус и ошибки & Пустой список документов --- корректный пустой индекс; чужой и несуществующий файл не открываются; запись в несуществующий каталог --- ошибка. \\
\hline
8 & Утилиты \texttt{indexer} и \texttt{lookup} & Индекс через командную строку, поиск терма и документа через \texttt{lookup}, коды возврата при ошибках во входных данных. \\
\hline
\end{tabular}
\end{center}

\section{Результаты}
\subsection{Автотесты}
Пройдены все 8 тестов. Журнал прогона (\texttt{results/tests.log}):

{\scriptsize
\VerbatimInput[breaklines, breakanywhere]{../results/tests.log}
}

\subsection{Индекс корпуса}
Индекс всего корпуса (\texttt{results/stats.tsv}, минимум из 3 запусков):
\begin{center}
\small
\begin{tabular}{|p{0.62\textwidth}|r|}
\hline
Показатель & Значение \\
\hline
Документов & 1\,000\,000 \\
\hline
Объём текста, МиБ & 5\,703,2 \\
\hline
Токенов & 995\,461\,170 \\
\hline
Термов & 3\,974\,421 \\
\hline
Записей в списках документов (пар терм--документ) & 363\,958\,293 \\
\hline
Средняя длина токена, символов & 4,83 \\
\hline
Средняя длина терма, символов & 8,03 \\
\hline
Размер индекса, МиБ & 1\,583,9 \\
\hline
Время индексации, с & 147,55 \\
\hline
\quad чтение файлов & 72,12 \\
\hline
\quad токенизация и вставка в словарь & 70,67 \\
\hline
\quad сортировка словаря & 1,87 \\
\hline
\quad запись индекса & 2,60 \\
\hline
Время на документ, мс & 0,147 \\
\hline
Время на килобайт текста, мкс & 25,3 \\
\hline
Скорость, МиБ/с & 38,7 \\
\hline
Пиковая память, МиБ & 2\,183,8 \\
\hline
\end{tabular}

\end{center}

Размер частей индекса:
\begin{center}
\small
\begin{tabular}{|l|r|r|}
\hline
Раздел & Байт & \% \\
\hline
Заголовок & 64 & 0,0 \\
\hline
Таблица документов & 12\,000\,000 & 0,7 \\
\hline
Таблица термов & 63\,590\,736 & 3,8 \\
\hline
Списки документов & 1\,455\,833\,172 & 87,7 \\
\hline
Пул строк & 129\,427\,722 & 7,8 \\
\hline
\end{tabular}

\end{center}
Почти девять десятых индекса (88\,\%) --- списки документов; их уменьшает сжатие (ЛР11: индекс 615~МиБ вместо 1\,584).

Пример работы \texttt{lookup}: \texttt{lviv} --- df 1\,508, \texttt{москва} --- df 40, \texttt{museum} --- df 124\,133, \texttt{the} --- df 987\,676 (из 1\,000\,000 документов), \texttt{zzzzqqq} --- df 0.

\subsection{Средняя длина терма и токена}
Средняя длина терма --- 8,03 символа, токена --- 4,83 (как в ЛР3). Средняя длина токена считается по всему тексту, т.\,е. каждое слово входит в неё столько раз, сколько встречается, а самые частые слова короткие: \texttt{the}, \texttt{of}, \texttt{and}, \texttt{in}, \texttt{a} (обратная зависимость частоты слова и его длины, закон Ципфа). В средней длине терма каждое слово учитывается один раз, а словарь в основном состоит из редких слов: 51\,\% термов встречаются в корпусе один раз (ЛР5) --- длинные редкие слова, имена собственные, числа (ISBN, телефоны), куски ссылок.

\subsection{Скорость индексации}
Весь корпус (5,6~ГиБ текста) индексируется за 148~с: 0,15~мс на документ, 25~мкс на килобайт текста, 38,7~МиБ/с. Из них 72~с --- чтение миллиона файлов, 71~с --- токенизация и вставка в словарь, 1,9~с --- сортировка словаря, 2,6~с --- запись файла. Пиковая память --- 2,1~ГиБ.

Индексация частей корпуса (\texttt{results/scale.tsv}; документы источников чередуются, пока источник не кончится):
\begin{center}
\footnotesize
\setlength{\tabcolsep}{4pt}
\begin{tabular}{|l|r|r|r|r|r|r|r|}
\hline
Док. & Текст, МиБ & Термов & Записей & Индекс, МиБ & Время, с & мкс/КиБ & Пам., МиБ \\
\hline
125\,000 & 819,8 & 1\,145\,386 & 50\,684\,675 & 234,6 & 18,05 & 21,5 & 802,2 \\
\hline
250\,000 & 1\,647,1 & 1\,788\,740 & 101\,873\,649 & 458,8 & 31,13 & 18,5 & 802,2 \\
\hline
500\,000 & 2\,837,0 & 2\,528\,687 & 182\,570\,787 & 809,4 & 65,69 & 22,6 & 1\,158,0 \\
\hline
1\,000\,000 & 5\,703,2 & 3\,974\,421 & 363\,958\,293 & 1\,583,9 & 147,55 & 25,3 & 2\,183,8 \\
\hline
\end{tabular}

\end{center}
При росте корпуса в 8 раз (по объёму текста --- в 7 раз) число записей в списках растёт почти пропорционально объёму (в 7,2 раза), а словарь --- в 3,5 раза: по закону Хипса $V \sim T^{\beta}$ с $\beta \approx 0{,}64$. Удельное время --- 18,5--25,3~мкс/КиБ, на всём корпусе больше всего. Причины две: миллион файлов хуже держится в кэше ОС, и словарь (4 млн записей по 32 байта и хеш-таблица на 8 млн ячеек --- около 160~МиБ) давно не помещается в кэш процессора, так что почти каждое обращение к нему --- промах кэша. Сортировка словаря растёт как $V \log V$, но остаётся малой (1,9~с из 148).

\subsection{Оптимальность и масштабирование}
Индексация не оптимальна. Её ограничивают:
\begin{itemize}
    \item чтение миллиона мелких файлов (около половины времени, 72 из 148~с): системные вызовы на открытие и чтение каждого файла и диск, когда файла нет в кэше ОС (см. ЛР3);
    \item случайный доступ к памяти при поиске в хеш-таблице и дописывании в списки (вторая половина времени): процессор ждёт память, а не считает;
    \item однопоточность: из 16 ядер занято одно.
\end{itemize}
Что можно ускорить:
\begin{itemize}
    \item хранить корпус в нескольких больших файлах или читать через \texttt{mmap}; читать файлы заранее в отдельном потоке;
    \item индексировать части корпуса в нескольких потоках, каждый со своим словарём, и сливать отсортированные части в конце;
    \item хранить хеш и короткий префикс терма прямо в ячейке таблицы, чтобы реже обращаться к записи терма; выделять списки документов из общего пула блоками, а не отдельным \texttt{realloc} на каждый терм; хранить списки в памяти сразу сжатыми (ЛР11);
    \item не копировать токен в \texttt{std::string} в токенизаторе, а передавать указатель на буфер.
\end{itemize}

Что будет при росте объёма (оценка по таблице и закону Хипса; сейчас записей в списках --- 364 млн, термов --- 3,97 млн, память --- 2,1~ГиБ, индекс --- 1,6~ГиБ):
\begin{itemize}
    \item \textbf{В 10 раз} (10 млн документов, 56~ГиБ текста): записей около 3,6 млрд, термов около 17 млн. Только списки документов в памяти --- 14--28~ГиБ (с запасом ёмкости динамических массивов): на этой машине (31~ГиБ) на грани, на обычной не поместится, нужна индексация блоками. Индекс --- около 16~ГиБ без сжатия и около 5~ГиБ со сжатием ЛР11. Текст уже не поместится в кэш ОС, время упрётся в диск --- десятки минут.
    \item \textbf{В 100 раз} (100 млн документов, 560~ГиБ текста): записей около 36 млрд --- около 150~ГБ, в память не помещаются. Нужна индексация блоками (SPIMI): строить индекс части корпуса, пока хватает памяти, записывать отсортированный блок на диск и в конце сливать блоки (здесь и пригодится слияние). Пул строк (заголовки и URL --- около 9~ГБ) превысит предел 4~ГиБ 4-байтовых смещений --- нужна новая версия формата с 8-байтовыми смещениями. Время --- часы, нужна параллельная индексация.
    \item \textbf{В 1000 раз} (1 млрд документов, 5,6~ТиБ текста): записей около 360 млрд, индекс около 1,5~ТБ без сжатия и около 0,4~ТБ со сжатием --- больше диска этой машины. Индекс нужно делить по документам между несколькими машинами и искать на всех параллельно.
\end{itemize}
Номера документов (4 байта) хватит до 4,29 млрд документов.

\section{Выводы}
Построен булев индекс корпуса из 1\,000\,000 документов: обратный индекс (3\,974\,421 терм в нижнем регистре, 364 млн записей \enquote{терм--документ}) и прямой индекс (заголовки, ссылки, число токенов) в одном бинарном файле собственного формата размером 1,6~ГиБ. Формат описан побайтово, рассчитан на расширение (версия, размеры заголовка и записей, флаги, смещения разделов --- им воспользовались ЛР4 и ЛР11), читается утилитой \texttt{lookup} и независимым читателем в тестах. Индексация занимает 148~с (0,15~мс на документ, 25~мкс/КиБ) и 2,1~ГиБ памяти, средняя длина терма --- 8,03 символа против 4,83 у токена.

Недостатки: индекс строится целиком в памяти и уже при росте корпуса в 10 раз не поместится в неё (нужна индексация блоками со слиянием); индексатор однопоточный и тратит половину времени на чтение мелких файлов; без сжатия списки документов занимают 88\,\% индекса (сжатие сделано в ЛР11); в словарь попадают мусорные термы токенизатора (числа, куски ссылок, склеенные слова).
